\section{Results}
\label{sec:results}

This section states the paper's empirical results in one place.
Construction details are in
Sections~\ref{sec:data-endpoints}--\ref{sec:robustness}.
Headline numerical results are generated from or hash-bound to frozen evidence.

\subsection{Headline branch separation}

The phenomenological low-redshift distance-sector endpoint~R
($H_{0}=\DistanceHZero$\kmsMpc{}; not presented as a complete cosmological model) differs substantially from both
Planck-compatible references C1 ($H_{0}=\CmbOneHZero$\kmsMpc{}) and
C2 ($H_{0}=\CmbTwoHZero$\kmsMpc{}, secondary sensitivity reference) across the constrained domain
$z\le1.8$.
The two R--CMB difference fields
$\Delta_{1}$ and $\Delta_{2}$ (defined in Section~\ref{sec:endpoint-repr})
share one dominant common
mode (common-mode fraction $\ge\GeometryCommonFractionMin{}$ per
interval), while the CMB--CMB difference
$\Delta_{C}$ is much smaller throughout: across the constrained bins,
the C1--C2 RMS difference is about 1--8\% of the R--CMB difference
(RMS ratios \GeometryCmbRatioLow{} and
\GeometryCmbRatioHigh{} on the constrained and high intervals).

\subsection{Compact reconstruction}

A five-coordinate order-3 (m3) map gives a compact empirical
description of the constrained-domain ($z\le1.8$) branch difference with local projected Jacobian rank~5
of~5 fitted coordinates ($\chi^{2}=\BridgeChiTwo{}$, only \BridgeDeltaCtrl{} above the
21-coordinate control; $\mathrm{AIC}\approx\BridgeAIC{}$ against
$\approx\BridgeAICCtrl{}$ for the control, reported descriptively as
parsimony rather than evidence). Selection is
predictive (held-out supernova and DESI $z\le1.8$ folds,
Section~\ref{sec:bridge}) and carries no significance of
its own: the matched constrained-domain null bootstrap tests representation complexity, not
existence of the endpoint separation (null bootstrap
$P(\mathrm{select}\ge m3)=\CvConstrBootProb{}$).
Zero ruler-clock split is allowed
($\Delta\chi^{2}(\epsilon{=}0)=\BridgeEpsDelta{}$).
The fitted difference reaches its largest
amplitude near $z\approx\GeometryPeakZ{}$; because this region also
carries substantial observational leverage, we do not interpret the
peak redshift as a fundamental physical scale
(Section~\ref{sec:peak-support}).

\subsection{Reference-branch robustness}

The frozen m3 basis transfers across the two Planck references with small
residual degradation (fraction \GeometryTransferDegradation{}):
parameters learned on
one reference construction predict the other separation. The two R--CMB fields
share one dominant direction
(cosine~\GeometryCosineSimilarity{}, principal angle
\GeometryPrincipalAngle$^{\circ}$). This reference-sensitivity
check shows the deformation is insensitive to the tested
CMB-reference construction; it does not assign a frequentist
p-value to the similarity (Section~\ref{sec:robustness-stats}).
Neither branch fits everything and no tension resolution is claimed.
High-redshift continuation through the
Lyman-alpha anchor remains unvalidated.
